\documentclass[letter,11pt]{article}
\usepackage[breaklinks]{hyperref}
\hypersetup{
    bookmarks=true,         % show bookmarks bar?
    unicode=false,          % non-Latin characters in Acrobat’s bookmarks
    pdftoolbar=true,        % show Acrobat’s toolbar?
    pdfmenubar=true,        % show Acrobat’s menu?
    pdffitwindow=false,     % window fit to page when opened
    pdfstartview={XYZ null null 1.00},    % disable zoom
    pdftitle={Advanced Malware Analysis Syllabus},    % title
    pdfauthor={Richard Zak},     % author
    pdfsubject={UMBC CYBER 691 Advanced Malware Analysis},   % subject of the document
    pdfkeywords={Malware, Data Science, Cybersecurity, Machine Learning}, % list of keywords
    pdfnewwindow=true,      % links in new PDF window
    colorlinks=false,       % false: boxed links; true: colored links
    linkcolor=red,          % color of internal links (change box color with linkbordercolor)
    citecolor=green,        % color of links to bibliography
    filecolor=magenta,      % color of file links
    urlcolor=cyan           % color of external links
}
\usepackage{graphicx}
\usepackage{fancyhdr}
\usepackage{multicol}
\pagestyle{fancy}
\usepackage[letterpaper, margin=1in]{geometry}
\geometry{letterpaper}
\usepackage{parskip} % Disable initial indent
\usepackage{color,soul} % Highligher

\usepackage[utf8]{inputenc}
\fancyhf{}
\renewcommand{\headrulewidth}{0pt} % Remove default underline from header package
\rhead{CYBR/CMSC 691 - Advanced Malware Analysis}
%\rhead{}
\lhead{\begin{picture}(0,0) \put(0,-10){\includegraphics[width=1.1cm]{Images/UMBC-vertical}} \end{picture}}
\cfoot{\thepage}
\rfoot{Spring 2020}
\lfoot{Advanced Malware Analysis}
\AtEndDocument{\vfill \hfill \LaTeX}

\begin{document}

\textbf{Semester:} Spring 2020\\
\textbf{Location:} ITE 241 \\
\textbf{Time:} Tuesdays 19:10 - 21:40 \\
\textbf{Instructor:} Richard Zak \\
\textbf{Email:} \href{mailto:richard.zak@umbc.edu?Subject=CYBR691}{richard.zak@umbc.edu}

\section*{Day/Hours Available}
\begin{itemize}
\item In Person: I can meet 30 minutes prior to class if you coordinate with me prior to the day of class. I can also stay after class for about 30 minutes.
\item In Email: Anytime, and I will respond within 48 hours.
\end{itemize}

\section*{Course Prerequisites}
\paragraph{}This is a BYOD course. Students must provide and bring to class a laptop capable of running Virtual Machines. Chromebooks, Android tablets, iPads, and similar devices are not acceptable. Please bring a system with at least 8 GB of RAM and 75-100 GB of available Hard Disk Space, 16GB of RAM or more is highly recommended.

\paragraph{}Completion of CYBR 620 or CMSC equivalent. Students must be proficient with Python, and preferably have had some exposure to assembly language. We will be using virtual machines with VirtualBox \url{https://www.virtualbox.org/}.

\section*{Course Description}
\paragraph{}This hands-on course will demonstrate how to build machine learning models for Windows executable malware. We will discuss reverse engineering, traditional malware analysis, static vs dynamic analysis, dataset curation, and available tools for analysis.

\section*{Overall Course Objectives}
\begin{multicols}{2}
\begin{itemize}
    \item Understand at a high level basic malware analysis techniques
    \item Understand the pros and cons of static analysis vs. dynamic analysis
    \item Understand the complexity of dealing with executable formats
    \item Understand the difference between domain-specific features vs. non domain-specific features
    \item Understand at a high level executable file formats
    \item Basic malware safety concepts, including use of virtual machines
    \item Basic machine learning concepts, including algorithms \& dataset curation
    \item Training \& evaluating machine learning models, and conceptual drift
\end{itemize}
\end{multicols}

\section*{Course Textbook}
\textbf{Malware Data Science} by Joshua Saxe \& Hillary Sanders, Publisher: No Starch Press; ISBN 978-1593278595, \url{https://nostarch.com/malwaredatascience}.

\subsection*{Optional Books}
These books are not part of the course, and are not required. Content from these books will not be on any exam. However, they may be helpful in your cyber data science career.
\begin{itemize}
    \item \textbf{Practical Malware Analysis} by Michael Sikorski and Andrew Honig, ISBN 978-1593272906. This book is highly regard in the malware analysis world. \url{https://nostarch.com/malware}
    \item \textbf{Practical Binary Analysis} by Dennis Andriesse, ISBN 978-1593279127. This book is a deep dive into crafting your own tools for analysis of executable file formats. \\ \url{https://nostarch.com/binaryanalysis}
    \item \textbf{Building Virtual Machine Labs} by Tony Robinson, ISBN 978-1546932635. This book might be useful if you want to learn more about setting up virtual machines to learn more about firewalls, routers, detonating malware, etc. \\ \url{https://www.amazon.com/Building-Virtual-Machine-Labs-Hands/dp/1546932631} \\ \url{https://leanpub.com/avatar}
\end{itemize}

\subsection*{Online Materials}
\paragraph{}Additional materials, including readings and the virtual machine, are available on a shared Google Drive folder located at \url{https://drive.google.com/open?id=1OMckkIhWibBllTvKLwngIrxY0zY_IS7B}.

\paragraph{VM} Two virtual machines for Virtual Box are available at the link above. They are the same, except one does not have a graphical interface, to reduce the size and resource requirements. Try the larger VM first, falling back to the other if needed.

\paragraph{Papers} Additional readings, which are relevant academic research papers, are available in the Google Drive link above. They will show you some of the cutting edge research being done in this field.

\section*{Course Work}
\paragraph{Labs}Most classes will conclude with a lab in order to practice the topic discussed in a technical, hands on fashion. The labs will begin as classwork, and become homework if not completed before the end of class. They are due before the following class. Make sure that your lab report includes the steps taken to complete the assignment, and includes a few screenshots or charts/graphs, as necessary.

\paragraph{Exams}There will be one exam, the final. It is designed to reinforce the topics discussed throughout the semester in order to promote retention of the information.

\paragraph{Projects}There will be two projects in which a paper or PowerPoint slides will be delivered. Regardless of paper vs. presentation, the content will be presented to the class. These projects are individual assignments, though discussing topics, ideas, \& concepts with your classmates is encouraged. You may swap the projects, so Project 2 may be submitted as Project 1, and so on. If you decide to do this, make sure it's clearly indicated.

\section*{Course Policies}
\subsection*{Grading}

\begin{center}
\begin{tabular}{ c c || c c}
\multicolumn{2}{c||}{\underline{Scale}} & \underline{Course Work} & \underline{Grade Distribution} \\
90\%-100\% & A & Final Exam & 25\% \\
80\%-89\% & B & Hands on Labs & 35\% \\
70\%-79\% & C & Project 1 & 20\% \\ 
60\%-69\% & D & Project 2 & 20\% \\
$<$60\%     & F & 
\end{tabular}
\end{center}

\section*{Writing and Research}
\paragraph{}Effective writing is critical to the intellectual life of university students and graduates within the workplace. Effective managers are usually effective communicators. Your work in this course must demonstrate your ability to master and effectively communicate course content. \\
Effective writing:
\begin{itemize}
    \item Meets the needs of the reader.
    \item Adequately covers the subject.
    \item Uses expected conventions of format and organization.
    \item Demonstrates use of credible reasoning and evidence.
    \item Satisfies standards of style and grammatical correctness.
    \item Requires 100\% compliance with UMBC's zero-tolerance policy regarding plagiarism.
\end{itemize}

\subsection*{Academic Integrity}
\paragraph{}By enrolling is this course, each student assumes the responsibilities of an active participant in UMBC's scholarly community in which everyone's academic work and behavior are held to the highest standards of honesty. Cheating, fabrication, plagiarism, and helping others to commit these acts are all forms of academic dishonesty, and they are wrong. Academic misconduct could result in disciplinary action that may include, but is not limited to, suspension or dismissal. To find useful information about avoiding plagiarism infractions through appropriate citations, or to read the full policy regarding student academic misconduct for the graduate school, please see \url{http://www.umbc.edu/provost/integrity}.

\section*{Schedule}
\textit{Note: The schedule is subject to change. However, any changes will be thoroughly discussed and well disseminated.}

\small
\begin{tabular}{c c l l c}
Date & Week & ~~~~~~~~~~~~Topic & ~~~~~~~~~~~~Lab & Chapter \\
\hline
Jan 27 & 1 & Introduction, Executable Formats & Lab Environment Overview & 1  \\
Feb 3 & 2 & Static \& Dynamic Analysis & Dynamic Analysis & 2 \& 3 \\
Feb 10 & 3 & Malware Networks, Shared Code Analysis & Malware Networks & 4 \& 5 \\
Feb 17 & 4 & Creating Machine Learning Malware Detectors & Feature Engineering & 6 \\
Feb 24 & 5 & Malware Neural Networks \& Model Evaluation & Model Evaluation & 7 \& 8 \\
Mar 2 & 6 & Cyber Security, Data Science, ML Overview  &  &  \\
Mar 9 & 7 & \hl{Project 1 Presentations}, Introduction to Go & & \\
Mar 16 & & Spring Break & & \\
Mar 23 & 8 & Guest Speaker &  &  \\
Mar 30 & 9 & More Malware Neural Networks & Explaining Feature Impact & 10 \& 11 \\
Apr 6 & 10 & Scalable $n$-gram Discovery & Building $n$-gram Models & \\
Apr 13 & 11 & Building a Logistic Regression Model & Regularization as Feature Selection & \\
Apr 20 & 12 & PDF \& RTF Malware Analysis & PDF Malware Model & \\
Apr 27 & 13 & Automated Yara Rule Generation & Automatic Yara Rule Creation & \\
May 4 & 14 & \hl{Project 2 Presentations}, Exam Review & & \\
May 11 & 15 & \hl{Final Exam} & \\
\end{tabular}

\small
\begin{tabular}{c c}
Week & Additional Reading \\
\hline
4 & investigation\_byte\_ngrams.pdf \& what\_can\_ngrams\_learn.pdf \\
5 & Variant\_Malware\_Testing\_Framework.pdf \\
9 & pe\_header\_ML.pdf \& malconv.pdf \\
10 & kilograms.pdf \\
15 & Malware\_Feature\_Selection\_Reinforcement\_Learning.pdf
\end{tabular}

\section*{Title IX Statement}
\paragraph{}As an instructor, I am considered a Responsible Employee, per UMBC’s Policy on Prohibited Sexual Misconduct, Interpersonal Violence, and Other Related Misconduct \footnote{\url{http://humanrelations.umbc.edu/sexual-misconduct/umbc-resource-page-for-sexual-misconduct-and-other-related-misconduct/}}. While my goal is for you to be able to share information related to your life experiences through discussion and written work, I want to be transparent that as a Responsible Employee I am required to report disclosures of sexual assault, domestic violence, relationship violence, stalking, and/or gender-based harassment to the University’s Title IX Coordinator.

\paragraph{}As an instructor, I also have a mandatory obligation to report disclosures of or suspected instances of child abuse or neglect\footnote{\url{http://www.usmh.usmd.edu/regents/bylaws/SectionVI/VI150.pdf}}.

\paragraph{}The purpose of these reporting requirements is for the University to inform you of options, supports and resources; you will not be forced to file a report with the police. Further, you are able to receive supports and resources, even if you choose to not want any action taken. Please note that in certain situations, based on the nature of the disclosure, the University may need to take action.
\end{document}
